%% posterstyle.tex: page, fonts, colours, blocks and helpers. Input from poster.tex after
%% venues.tex has set \posterwidth, \posterheight and \posterscale.
%% The rules behind these choices are in GUIDELINE.md, section 14.

\documentclass[final,xcolor=table]{beamer}
\usepackage[size=custom,width=\posterwidth,height=\posterheight,scale=\posterscale]{beamerposter}

%% Fonts: Helvetica for all text (high contrast and a large x-height at a distance), with
%% Times-style maths as in the paper. Latin Modern looks thin and washed out at poster scale.
\usepackage[T1]{fontenc}
\usefonttheme{professionalfonts}   % keep the newtx maths; beamer would otherwise reset it
\usepackage{amsmath}
\usepackage{newtxmath}
\usepackage{iftex}
\ifPDFTeX
  \usepackage[scaled=0.95]{helvet}
\else  % XeLaTeX / LuaLaTeX
  \usepackage{fontspec}
  \setsansfont{texgyreheros}[Extension=.otf, UprightFont=*-regular, BoldFont=*-bold,
    ItalicFont=*-italic, BoldItalicFont=*-bolditalic, Scale=0.95]
\fi
\renewcommand{\familydefault}{\sfdefault}
\usepackage{graphicx,booktabs,tabularx,tikz,pgfplots,qrcode,fontawesome5}
\usepackage[savepos]{zref}
\usepackage[most]{tcolorbox}
\usetikzlibrary{arrows.meta,positioning,calc}
\usepgfplotslibrary{groupplots}
\pgfplotsset{compat=1.18}

\setbeamertemplate{navigation symbols}{}
\setbeamertemplate{headline}{}
\setbeamertemplate{footline}{}
\setbeamersize{text margin left=0cm,text margin right=0cm}
\setbeamercolor{normal text}{fg=black}
\usebeamercolor[fg]{normal text}

%% ---------------- colours ----------------
%% One saturated accent for the first title line, block headings, number discs and rules.
%% Keep it away from the plot roles: burnt orange headings with a blue "ours" work well.
\providecommand{\accenthex}{BF5700}
\definecolor{accent}{HTML}{\accenthex}
\definecolor{panel}{HTML}{EFEFEF}       % flat grey panels: cards, equation boxes, chips
\definecolor{rulegray}{HTML}{ABABAB}    % header rule
\definecolor{headgray}{RGB}{51,63,72}   % affiliations, small header labels
\definecolor{goodgreen}{RGB}{27,138,90}
\definecolor{badred}{RGB}{200,50,40}

%% posterplots.tex: the house figure style (GUIDELINE.md) for pgfplots at poster scale.
%% Input after pgfplots is loaded (posterstyle.tex does this). Reusable as is.
%%
%% Scale: the guide's sizes assume 8 pt axis labels. Line widths, markers and gaps below are
%% the guide's values times about 3.2; text uses the poster sizes set in \figtick etc.
%%   - sans text (the guide's setting for posters), lowercase axis labels with units
%%   - thin full frame, no tick marks, light grid behind the data
%%   - fixed colour roles, each with a second encoding (marker, dash or hatch)
%%   - bars in light shades mixed with white (46%, ours 66%, oracle 30%), baselines hatched
%%   - one legend, ours first; text in ink, never in a series colour

\usetikzlibrary{patterns.meta}

%% ---------------- colour roles (GUIDELINE.md, section 4) ----------------
\definecolor{oursc}{HTML}{2E7BD6}     % ours: always
\definecolor{ac}{HTML}{C9302C}        % a: strongest baseline
\definecolor{bc}{HTML}{BF8700}        % b: baseline
\definecolor{cc}{HTML}{A23BA6}        % c: baseline
\definecolor{dc}{HTML}{159E86}        % d: baseline
\definecolor{basec}{HTML}{8C8C8C}     % base: untreated model, background points, connectors
\definecolor{oraclec}{HTML}{333333}   % oracle: upper bound or reference
\colorlet{variantc}{oursc!60!white}   % a variant of ours shares its hue
\definecolor{inkc}{HTML}{1A1A1A}
\definecolor{mutedc}{HTML}{666666}
\definecolor{gridc}{HTML}{E4E4E4}

%% Figure text sizes, captured here so every plot uses the same sizes even inside a column
%% that enlarges its body text with \largertext. At scale 1.24: 25.7 pt, 30.9 pt, 21.4 pt.
\let\figtick\small          % ticks, legends
\let\figlabel\normalsize    % axis labels, panel titles
\let\figvalue\footnotesize  % value labels

%% ---------------- series styles (TikZ, so they also work in legends) ----------------
\tikzset{
  s ours/.style={oursc, line width=5.2pt, mark=*, mark size=5.6pt, mark options={solid}},
  s a/.style={ac, line width=3.8pt, mark=square*, mark size=4.6pt, mark options={solid}},
  s b/.style={bc, line width=3.8pt, dash pattern=on 19pt off 5.7pt,
    mark=triangle*, mark size=6.2pt, mark options={solid}},
  s c/.style={cc, line width=3.8pt, dash pattern=on 11.5pt off 3.8pt on 3.8pt off 3.8pt,
    mark=diamond*, mark size=6.2pt, mark options={solid}},
  s d/.style={dc, line width=3.8pt, dash pattern=on 3.8pt off 3.8pt,
    mark=triangle*, mark size=6.2pt, mark options={solid, rotate=180}},
  s base/.style={basec, line width=3.2pt, mark=*, mark size=4.6pt, mark options={solid}},
  s oracle/.style={oraclec, line width=3.2pt, dash pattern=on 12.8pt off 6.4pt, mark=none},
  s variant/.style={variantc, line width=3.8pt, dash pattern=on 19pt off 5.7pt, mark=none},
  % bars: light shade mixed with white; baselines keep a hatch in the full role colour
  hatch/.style 2 args={postaction={pattern={Lines[angle=#2, distance=0.34cm, line width=1.6pt]}, pattern color=#1}},
  b ours/.style={fill=oursc!66!white, draw=none},
  b a/.style={fill=ac!46!white, draw=none, hatch={ac}{45}},
  b b/.style={fill=bc!46!white, draw=none, hatch={bc}{-45}},
  b c/.style={fill=cc!46!white, draw=none,
    postaction={pattern={Dots[distance=0.34cm, radius=1.6pt]}, pattern color=cc}},
  b d/.style={fill=dc!46!white, draw=none,
    postaction={pattern={Hatch[distance=0.34cm, line width=1.6pt]}, pattern color=dc}},
  b base/.style={fill=basec!46!white, draw=none},
  b oracle/.style={fill=oraclec!30!white, draw=none},   % reference bars lighter, so ours reads first
  b variant/.style={fill=oursc!40!white, draw=none},
  % dumbbell dots: hollow = before (e.g. peak), filled = after (e.g. final)
  d before/.style={draw=#1, fill=white, line width=3.4pt},
  d after/.style={draw=#1, fill=#1, line width=3.4pt},
}

%% ---------------- legends and keys (set in the figure font) ----------------
% \legline{style}{label}  line sample      \legbar{style}{label}  bar swatch
% \legdot{style}{label}   dot              \figlegend{entries}    one centred row under a figure
\newcommand{\legline}[2]{\mbox{\tikz[baseline=-0.6ex]\draw[#1, mark indices={2}] plot coordinates {(0,0) (1.1cm,0) (2.2cm,0)};\hspace{0.4em}#2}}
\newcommand{\legbar}[2]{\mbox{\tikz[baseline=-0.1ex]\path[#1] (0,-0.3) rectangle (1.0,0.5);\hspace{0.4em}#2}}
\newcommand{\legdot}[2]{\mbox{\tikz[baseline=-0.6ex]\path[#1] (0,0) circle (8pt);\hspace{0.4em}#2}}
\newcommand{\figlegend}[1]{{\centering\sffamily\figtick\color{inkc}#1\par}}
% A text line that starts with a series' line sample. Lines like this under a plot can be its
% legend: \keyline{s a}{\textbf{Baseline A:} what happens to it.}
\newcommand{\legkey}[1]{\tikz[baseline=-0.6ex]\draw[#1, mark indices={2}] plot coordinates {(0,0) (0.9cm,0) (1.8cm,0)};}
\newcommand{\keyline}[2]{{\hangindent=2.5cm\hangafter=1\noindent\makebox[2.5cm][l]{\legkey{#1}}#2\par}\vspace{6pt}}

%% ---------------- axis style ----------------
\pgfplotsset{
  posteraxis/.style={
    scale only axis,
    axis line style={draw=inkc, line width=1.9pt},
    xmajorgrids, ymajorgrids, grid style={gridc, line width=1.6pt},
    tick style={draw=none},
    /pgf/number format/assume math mode=true,   % sans digits
    tick label style={font=\sffamily\figtick, inkc},
    label style={font=\sffamily\figlabel, inkc},
    title style={font=\sffamily\figlabel, inkc, yshift=-0.2cm},
    every node near coord/.append style={font=\sffamily\figvalue, inkc},
  },
  % bar charts: no grid along the categorical axis
  posterbars/.style={posteraxis, xmajorgrids=false},
}

%% Second-level category label under a group of ticks, e.g. the model family under its sizes.
%% Use inside an axis with clip=false: \famlabel{symbolic x coord}{text}
\newcommand{\famlabel}[2]{\node[font=\sffamily\figtick, text=inkc, anchor=north, yshift=-1.3cm] at (axis cs:#1,0) {#2};}

%% Dumbbell row inside an axis: \Dumbbell{colour}{row}{before}{after}
\newcommand{\Dumbbell}[4]{%
  \draw[basec, line width=5pt] (axis cs:#4,#2) -- (axis cs:#3,#2);
  \path[d before=#1] (axis cs:#3,#2) circle (9pt);
  \path[d after=#1]  (axis cs:#4,#2) circle (9pt);
}
   % colour roles, plot styles, \pu and the scaled line widths

%% Blue (the "ours" colour) is reserved for our method and our key results; titles never use it.
%% Text uses a slightly darker shade than the plot lines, so small bold text stays readable.
\colorlet{oursT}{oursc!82!black}

%% ---------------- blocks ----------------
%% \pblock{title}{body}: a white rounded box with a bold black outline. Inside: a number disc and
%% a bold heading in the accent colour, centred, over an accent rule. Blocks are numbered in
%% the order they are typeset, so a layout that drops a block renumbers the rest.
%% Headings are \Large and must fit on one line beside the disc: state the finding briefly.
\newcounter{pblocknum}
\newlength{\headW}
\pkdim{\blockrule}{5pt}
\pkdim{\headrule}{4pt}
\newcommand{\pblock}[2]{%
  \stepcounter{pblocknum}%
  \begin{tcolorbox}[colback=white, colframe=black, boxrule=\blockrule, arc=0.6\pu, boxsep=0pt,
      left=0.5\pu, right=0.5\pu, top=0.35\pu, bottom=0.4\pu, before skip=0pt, after skip=0pt]
  \settowidth{\headW}{\sffamily\bfseries\Large #1}%
  \ifdim\dimexpr\headW+2.5\pu\relax>\linewidth
    \PackageWarningNoLine{posterstyle}{The heading of block \arabic{pblocknum} wraps to a second line:
      shorten '#1'}%
  \fi
  {\centering\tikz[baseline=(n.base)]\node[circle, fill=accent, text=white, inner sep=0pt,
      minimum size=1.85\pu, font=\large\sffamily\bfseries](n){\arabic{pblocknum}};\hspace{0.45em}%
   {\sffamily\bfseries\Large\color{accent}#1}\par}
  \vspace{0.2\pu}{\color{accent}\rule{\linewidth}{\headrule}}\par\vspace{0.45\pu}
  {#2\par}
  \end{tcolorbox}}
%% space between stacked blocks: at least 0.8 cm (at scale 1.36), plus an even share of the
%% column's spare height
\newcommand{\blockgap}{\par\vspace{0.8\pu plus 1fill}}

%% ---------------- helpers ----------------
% one-line takeaway under a figure
\newcommand{\take}[1]{{\small #1\par}}
% small print (never smaller than \footnotesize)
\newcommand{\smallprint}[1]{{\footnotesize #1\par}}
% our method or one of our key results, in the "ours" colour
\newcommand{\ours}[1]{{\color{oursT}\bfseries #1}}
% another highlighted keyword in a given colour, e.g. a baseline's failure: \hl{ac}{collapses}
\newcommand{\hl}[2]{{\color{#1}\bfseries #2}}
% setup chip: \chip{\faIcon}{text}
\newcommand{\chip}[2]{\tcbox[on line, colback=panel, colframe=panel, arc=0.25\pu, boxsep=0pt,
  left=0.25\pu, right=0.25\pu, top=0.2\pu, bottom=0.2\pu]{\footnotesize{\color{accent}#1}\ #2}}
% grey reference row in a results table; \leavevmode keeps the colour from opening an empty
% line at the start of a tabularx X cell
\newcommand{\refcell}[1]{\leavevmode{\color{black!72}#1}}
\newcommand{\yes}{{\color{goodgreen}\faCheck}}
\newcommand{\no}{{\color{badred}\faTimes}}
\setbeamertemplate{itemize items}{\color{accent}\faCaretRight}
% equation slot of fixed height, so text under side-by-side boxes lines up
\providecommand{\eqslotH}{6.5\pu}
\newcommand{\eqslot}[1]{\par\parbox[c][\eqslotH][c]{\linewidth}{\centering\small #1}\par}
% takeaway at the end of the last block: one sentence, on a light panel in the "ours" hue.
% Limitations are said at the poster, not printed on it.
\newcommand{\takeaway}[1]{%
  \begin{tcolorbox}[colback=oursc!9, colframe=oursc!9, arc=0.45\pu, left=0.5\pu, right=0.5\pu,
      top=0.4\pu, bottom=0.4\pu, before skip=0pt, after skip=0pt]
    \raggedright\hyphenpenalty=10000\small\ours{Takeaway:} #1\par
  \end{tcolorbox}}
% A column with spare height can step every body size up by one. Call at the start of the
% column. Plot text is unaffected (its sizes are captured in posterplots.tex).
\newcommand{\largertext}{\let\origsmall\small \let\orignormalsize\normalsize
  \let\scriptsize\footnotesize \let\footnotesize\origsmall \let\small\orignormalsize}

% grey dashed box standing in for a figure or logo that is not there yet
\newcommand{\placeholder}[3]{% width, height, label
  \tikz\node[draw=black!35, dashed, line width=0.07\pu, rounded corners=0.28\pu, minimum width=#1,
    minimum height=#2, align=center, font=\sffamily\small, text=black!50]{#3};}
% logo: the file if it exists, otherwise a placeholder of about the right size
% \logofile{height}{width of placeholder}{path with extension}{label}
% A logo never gets wider than the space it sits in.
\newcommand{\logofile}[4]{\IfFileExists{#3}{\includegraphics[height=#1, width=\linewidth, keepaspectratio]{#3}}%
  {\setlength{\cardW}{\minsize{#2}{\linewidth-0.15\pu}}\placeholder{\cardW}{#1}{#4}}}  % room for the dashed outline

% icon size for cards: 54 pt at scale 1.36
\pgfmathsetmacro{\iconsize}{39.7*\posterscale}
% \minsize{a}{b}: the smaller of two length expressions. Layouts use it so a figure is as tall
% as its width suggests, but never taller than its share of the column height.
\newcommand{\minsize}[2]{\ifdim\dimexpr#1\relax<\dimexpr#2\relax\dimexpr#1\relax\else\dimexpr#2\relax\fi}
% scratch lengths for sizes computed from the column width
\newlength{\cardW}\newlength{\cardT}\newlength{\figH}
